\subsection{Evolution of Python Type System}
The ability to omit explicit types in Python allows developers to write concise
code and rapidly prototype programs.
However, this high degree of flexibility can reduce code readability,
complicate maintenance, and increase the likelihood of unintended type-related
bugs.
To mitigate these issues, Type Annotations were introduced in Python 3.5 for
variables and in Python 3.6 for function parameters and return.
Type Annotations are optional and do not affect the runtime behavior of a
program, but they enable developers to explicitly state type information.
This improves readability and developer productivity while serving as a form of
contract that helps enhance program reliability.

Following the introduction of Type Annotations, various static type checkers
have been actively developed in the Python ecosystem.
Before this change, static analysis tools had to rely entirely on type
inference, which is notoriously difficult in dynamic languages and often leads
to imprecise analyses.
With explicit annotations, static analyzers can perform significantly more
accurate checks by leveraging developer-provided type information rather than
depending solely on inference.

Python thus adopts a Gradual Typing paradigm~\cite{gradualtyping} that combines
the advantages of both dynamic and static type systems.
Under Gradual Typing, parts of the program where types can be determined
statically are checked statically, while the rest continue to operate under the
dynamic type system.
This hybrid approach allows Python to achieve both flexibility and reliability
in large-scale software development.
The concept of Gradual Typing illustrates with the \autoref{lst:gradualTyping}

\begin{listing}[t]
\begin{minted}[linenos, numbers=left, frame=lines, framesep=2mm, xleftmargin=10pt, fontsize=\footnotesize]{python}
def foo(a: int, b: int):
    return a + b

# No Error
a = foo(1, 2) 

# TypeError only in static type checker
b = foo(2, 3.14)

# TypeError in both static type checker and runtime
c = foo(3, "hello")
\end{minted}
\caption{Example Python program demonstrating object compatibility determined by duck typing.}
\label{lst:gradualTyping}
\end{listing}

In this example, the function foo explicitly specifies that both a and b should
be of type int.
This example highlights the following three key observations:

\begin{itemize}
  \item foo(1, 2) satisfies the type specification, so neither the static type
    checker nor the runtime reports any errors.

  \item foo(2, 3.14) causes a type error in the static type checker because a
    float is passed where an int is expected. However, the Python runtime accepts
    the call and executes it, demonstrating that annotations are not enforced at
    runtime.

  \item foo(3, "hello") produces a type error in the static type checker, and
    additionally, it raises a runtime TypeError because Python cannot add an
    integer and a string.
\end{itemize}

This example illustrates how Python's Gradual Typing performs static checking
where type information is available, while maintaining dynamic typing behavior
where static analysis is insufficient.
This hybrid model provides the reliability of a static type system without
sacrificing the flexibility of dynamic typing, making it suitable for
large-scale software development.

